% !TEX program = xelatex

\documentclass{resume}
\usepackage{zh_CN-Adobefonts_external} % Simplified Chinese Support using external fonts (./fonts/zh_CN-Adobe/)
%\usepackage{zh_CN-Adobefonts_internal} % Simplified Chinese Support using system fonts

\begin{document}
\pagenumbering{gobble} % suppress displaying page number

\name{Li Wang 王力}

\basicInfo{
  \email{223040245@link.cuhk.edu.cn} \textperiodcentered\ 
  \faDesktop{\href{https://wwwwwli.github.io/}{~Li Wang Homepage}} \textperiodcentered\ 
  \faGraduationCap {\href{https://scholar.google.com/citations?user=NLQPDrsAAAAJ&hl}{~Google Scholar}} }

\small \textbf{Research Interests:} Safety, Audio Deepfake Detection, Audio Language Model


\section{Education}
\datedsubsection{\textbf{The Chinese University of Hong Kong , Shenzhen}}{Sep. 2023 --}
\begin{itemize}
    \item PhD in Computer Science
    \item Areas of Study: audio deepfake and anti-spoofing
    \item Advisor: \href{https://sds.cuhk.edu.cn/teacher/641}{Zhizheng Wu} and \href{https://sds.cuhk.edu.cn/teacher/498}{Haizhou Li}
\end{itemize}


\datedsubsection{\textbf{Peking University} (Exam-exempted postgraduate)}{Sep. 2019 -- Jun. 2022}
\begin{itemize}
    \item Master in Electronic and Computer Engineering (ECE)
    \item Areas of Study: spoken keyword spotting and speaker verification
    \item Advisor: \href{https://web.pkusz.edu.cn/adsp/}{Yuexian Zou}
\end{itemize}

\datedsubsection{\textbf{University of Electronic Science and Technology of China}}{Sep. 2015 -- Jun. 2019}
\begin{itemize}
    \item \textit{B.Eng.} in Information and Software Engineering
    \item GPA: 3.69/4.00, Ranking: 8/162(Top 4.9\%)
    \item Core courses: Random Signal Analysis (99), Digital Signal Processing (95)
\end{itemize}

\section{Publications}
\begin{enumerate}\itemsep 0.5em
  \item \textbf{Li Wang}, Xiao Lei, Haorui He, Lei Wang, Jie Shi, Zhizheng Wu. \href{https://ieeexplore.ieee.org/document/11334038}{[PDF]} \\Over-the-Air Adversarial Attacks and Detection for Automatic Speaker Verification. \textit{IEEE/ACM Transactions on Audio, Speech, and Language Processing (TASLP), 2026.}
  \item Peizhuo Liu*, \textbf{Li Wang*}, Renqiang He*, Haorui He, Lei Wang, Huadi Zheng, Jie Shi, Tong Xiao, Zhizheng Wu. \href{https://arxiv.org/abs/2409.11308}{[PDF]} \\SpMis: An Investigation of Synthetic Spoken Misinformation Detection. \textit{Proceedings of the IEEE Spoken Language Technology Workshop (SLT 2024).} (\textbf{Best Paper Finalist, Top 2.5\%})
  \item \textbf{Li Wang}, Jiaqi Li, Yuhao Luo, Jiahao Zheng, Lei Wang, Hao Li, Ke Xu, Chengfang Fang, Jie Shi, Zhizheng Wu. \href{https://arxiv.org/abs/2310.05369}{[PDF]} \\AdvSV: An Over-the-Air Adversarial Attack Dataset for Speaker Verification. \textit{Proceedings of the 49th IEEE International Conference on Acoustics, Speech and Signal Processing (ICASSP 2024).}
 \item Jiaqi Li, \textbf{Li Wang}, Liumeng Xue, Lei Wang, Zhizheng Wu. \href{https://arxiv.org/abs/2310.05354}{[PDF]} \\An Initial Investigation of Neural Replay Simulator for Over-the-Air Adversarial Perturbations to Automatic Speaker Verification. \textit{Proceedings of the 49th IEEE International Conference on Acoustics, Speech and Signal Processing (ICASSP 2024).}
  \item \textbf{Li Wang}, Rongzhi Gu, Weiji Zhuang, Peng Gao, Yujun Wang, Yuexian Zou. \href{https://ieeexplore.ieee.org/document/9747878}{[PDF]} \\Learning Decoupling Features Through Orthogonality Regularization. \textit{Proceedings of the 47th IEEE International Conference on Acoustics, Speech and Signal Processing (ICASSP 2022).}
  \item \textbf{Li Wang}, Rongzhi Gu, Nuo Chen, Yuexian Zou. \href{https://www.isca-speech.org/archive/interspeech_2021/wang21da_interspeech.html}{[PDF]} \\ Text Anchor Based Metric Learning for Small-Footprint Keyword Spotting. \textit{Proceedings of the 22nd Annual Conference of the International Speech Communication Association (Interspeech 2021).}

\end{enumerate}


\section{Experience}
\datedsubsection{\textbf{Shenzhen Research Institute of Big Data (SRIBD)}}{Nov. 2022 -- Aug. 2023}
{\small \role{Research Assistant, Institute of Basic Theory and Algorithm of Big Data, Shenzhen, China}{}}
Leading a project with Huawei on detecting physical attacks in speaker verification technology.
\begin{itemize}
    \item Successfully generated a dataset of 1.1 million data samples for physical attacks on speaker verification.
    \item Determined the technical route for the defense method, which was fully affirmed and supported by Huawei.
    \item Contributed to developing a reliable defense mechanism for speaker verification against physical attacks.
    \item Demonstrated strong leadership, problem-solving, and technical skills throughout the project.
\end{itemize}


\datedsubsection{\textbf{International Digital Economy Academy (IDEA)}}{Jul. 2022 -- Oct. 2022}
{\small \role{Financial Machine Learning Engineer, AI Finance and Deep Learning Research Center, Shenzhen, China}{}}
Create financial factors (features) and design deep learning models to predict stock returns.


\section{Internships}
\datedsubsection{\textbf{AILAB, ByteDance}}{Jul. 2021 -- Mar. 2022}
\begin{itemize}
    \item Topic: Audio fingerprinting and music segmentation
    \item Job Description: Proposed the audio fingerprint temporal localization algorithm. The accuracy rate was improved from 95.80\% to 97.90\%, and the IOU of temporal localization was 0.62.
\end{itemize}

\datedsubsection{\textbf{Audio Engineering Department, HuaWei MTI}}{Apr. 2021 -- Jun. 2022}
\begin{itemize}
    \item Topic: Music separation
    \item Job Description: Replicated D3Net and trained and tested on the MUSDB dataset, separating vocals, drums, bass and other
tracks separately, with a test average SDR of 6.43.
\end{itemize}

\datedsubsection{\textbf{AI Lab, Tencent}}{May. 2020 -- Nov. 2020}
\begin{itemize}
    \item Topic: Speech separation and speaker verification
    \item Job Description: Testing the performance of the speech separation model GALR on a speaker verification task, collating
speaker information from 2,000,000 pieces of audio data.
\end{itemize}

\section{Projects}
\datedsubsection{\textbf{Research on Voice Wake-up Technology of Artificial Intelligence}}{Jan. 2021 -- Jan. 2022}
\role{Leader}{Collaborative project between the lab and Xiaomi during graduate school.}
\begin{itemize}
  \item Communicated with Xiaomi to track project progress, understand data requirements.
  \item Optimized network design and data enhancement to improve model performance.
  \item Reduced the false rejection rate from 20\% to 2\% and achieved an accuracy rate of 99.7\%.
  \item Collaborated on personalized wake-up task and submitted ICASSP2022 paper as first author.
\end{itemize}


\section{Services}
\datedsubsection{\textbf{Reviewer}}{}
\begin{itemize}
  \item IEEE/ACM Transactions on Audio, Speech, and Language Processing (IF 4.1)
  \item IEEE International Conference on Acoustics, Speech and Signal Processing
  \item IEEE Signal Processing Letters (IF 3.201)
  \item National Conference on Man-Machine Speech Communication (NCMMSC)
  \item International Conference on Asian Language Processing (IALP)
  \item Computer Speech \& Language (IF 3.252)
  \item Workshop on Automatic speech Recognition and Understanding
\end{itemize}

\datedsubsection{\textbf{Standards Development}}{}
\begin{itemize}
  \item Contributed to the development of ``Method for Identifying AI-Generated Synthetic Content'' (Metadata Implicit Identification for Text Files), published by National Technical Committee 260 on Cybersecurity of SAC (TC260).
\end{itemize}

\datedsubsection{\textbf{Student Volunteer}}{}
\begin{itemize}
  \item IEEE Spoken Language Technology Workshop 2024
\end{itemize}

\datedsubsection{\textbf{Teaching Assistant}}{}
\begin{itemize}
  \item 2023 Fall, CSC1001 Introduction to Computer Science: Programming Methodology, CUHK-Shenzhen.
  \item 2024 Spring, CSC3160/AIR6063 Fundamentals of Speech and Language Processing/Spoken Language Processing, CUHK-Shenzhen.
  \item 2024 Fall, DDA3020 Machine Learning, CUHK-Shenzhen.
  \item 2025 Spring, DDA3020 Machine Learning, CUHK-Shenzhen.

\end{itemize}


\section{Honors and Awards}
\begin{enumerate}
  \item National Second Prize of National University Student Information Security Competition \textit{全国大学生信息安全竞赛国家级二等奖}(2018)
  \item Undergraduate People's First Class Scholarship \textit{本科生人民一等奖学金}(2015)
  \item Undergraduate People's First Class Scholarship \textit{本科生人民一等奖学金}(2016)
  \item Undergraduate People's First Class Scholarship \textit{本科生人民一等奖学金}(2017)
  \item Outstanding Volunteer of the 9th Paralympic Games of the People's Republic of China. \textit{中华人民共和国第九届残疾人运动会优秀志愿者}(2015)
  \item Peking University 4x100m Runner-up in the School Games. \textit{北京大学校运会4x100米亚军}(2021)
\end{enumerate}


\section{Skills}
\begin{itemize}
  \item Erhu Grade 10.
  \item Programming Languages: proficient in Python, familiar with Java, and know about C++.
  \item Deep Learning Frameworks: proficient in PyTorch.
\end{itemize}

%% Reference
%\newpage
%\bibliographystyle{IEEETran}
%\bibliography{mycite}
\end{document}
